\documentclass[10pt,a4paper]{amsart}
\usepackage[margin=19mm]{geometry}
\usepackage{amsmath,amssymb,mathtools}
\usepackage[scr=rsfs,cal=euler]{mathalfa}
\usepackage{xfrac}
\usepackage[unicode,hidelinks,bookmarks=false]{hyperref}
\newcommand{\ie}{i.e.\ }
\newcommand{\cf}{cf.\ }
\newcommand{\resp}{respectively\ }
\newcommand{\Subsection}{\subsection}
\providecommand{\nobreakdash}{\nobreak\hbox{-}}
\setlength{\emergencystretch}{2em}
\multlinegap0pt
\usepackage{amssymb}
\newtheorem{lemma}{Lemma}
\newcommand{\Si}{\Sigma}
\renewcommand{\a}{\alpha}
\newcommand{\bV}{\mathbb V}
\newcommand{\be}{\begin{equation}}
\newcommand{\bes}{\begin{equation*}}
\newcommand{\cC}{{\mathcal C}}
\newcommand{\cU}{{\mathcal U}}
\newcommand{\ee}{\end{equation}}
\newcommand{\ees}{\end{equation*}}
\newcommand{\g}{\gamma}
\renewcommand{\k}{\kappa}
\newcommand{\p}{\partial}
\renewcommand{\to}{\rightarrow}
\newcommand{\tr}{\mathrm{tr}}
\title{Well-Posed Geometric Boundary data in General Relativity, I: Dirichlet boundary data}
\author{Zhongshan An, Michael T. Anderson}
\date{Source: arXiv:2505.07128v2 (CC BY 4.0)}
\begin{document}
\maketitle

\noindent\emph{Source and licence.} Well-Posed Geometric Boundary data in General Relativity, I: Dirichlet boundary data. Version arXiv:2505.07128v2 (CC BY 4.0), distributed by arXiv under the Creative Commons Attribution 4.0 International licence, http://creativecommons.org/licenses/by/4.0/, as recorded in the arXiv licence metadata for this exact version and shown on the abstract page https://arxiv.org/abs/2505.07128v2. Full licence notice: \texttt{licenses/arxiv-2505.07128-CC-BY-4.0.txt}.\par\smallskip

\noindent\textit{Editorial scope.} Counted unit: the article's lemma Man1 (source0.tex lines 1008--1010). The article's complete original proof of it is delivered with it (source0.tex lines 1011--1054, one \texttt{proof} environment of 44 lines). No mathematics is written, repaired or re-derived: the statement and the proof are the article's own lines, in the article's own notation, and no retained line is substituted or re-flowed.  Retained uncounted as the source-local context the proof needs: lines 900--907.  Nothing was omitted from any retained span, so the package carries no omission record.  No retained line contains a citation, so the wrapper carries no bibliography and no bibliography entry.  No retained line contains a figure environment, an image-inclusion command or drawing code, so no image is delivered and the package declares no asset dependency.  Editorial formatting only, in addition: an AMS \texttt{amsart} preamble that restates the article's own declarations and macros used by the retained lines, namely the environments \texttt{lemma} on separate counters, together with the standard packages those lines need.  The retained environments print in this document as Lemma 1 in order of appearance, whereas the article prints the same items with the numbering of its own full document.  The complete native source of the article as distributed in the arXiv e-print for this exact version is bundled unchanged as \texttt{sources/177939/source0.tex}.
\begin{lemma}\label{alpha}
For any $C^1$ metric $g$, one has the relation 
\be \label{vel}
H_{\Si}^{\cC} = \sqrt{1+\a^2}\tr_\Si \k - \a H_\Si^S,
\ee
where $H_{\Si}^S$ and $H_{\Si}^{\cC}$ are the mean curvatures of $\Si$ in $(S, \g_S)$ and $(\cC, \g_{\cC})$ respectively. 
Consequently, under Assumption (*), $\a$ is determined by the data $(\g_S, \k, \g_{\cC})$. 
\end{lemma}

\begin{lemma}\label{Man1}
For $s$ sufficiently large (depending only on the dimension $n$), the space $\cU_c^1$ is a closed Hilbert submanifold of $\cU$. 
\end{lemma}

\begin{proof}
    Define the following map 
\be \label{CompOp1}
W^1: \cU \to Sym(\Si)^2 \times \bV'(\Si) \times \bV(\Si) \times \bV_{\cC}(\Si),
\ee
\bes
W^1(\g_S, \k, \nu_S, V_S, \g_{\cC}, V_{\cC}) = 
\left\{ 
\begin{array}{l}
(\g_S - \g_{\cC})|_{\Si} \\
\p_t \g_{\cC} - \chi(\p \g_S, \k, \p g_\cC) \\
\nu - \chi(\a) \\
(V_{\cC} - V_S)|_{\Si} \\
V_S^a - \chi^a(V_S^1,\p \g_S,\k,\p \nu,\p \g_\cC) \\
\end{array}\right.
\ees
In the above, all data on the right are evaluated at $\Si$; also $\bV_{\cC}(\Si)$ denotes vector fields at $\Si$ tangent to $\cC$. 
By definition, the inverse image $W^{-1}({\bf 0}) = \cU_c^1$. 

  We claim that $W^1$ is a submersion at $\cU_c^1$, i.e.~the derivative $DW^1$ is everywhere surjective on $\cU_c^1$. To see this, consider 
infinitesimal deformations of the form $(0, 0, \nu', V_S', \g_{\cC}', V_{\cC}' ) \in T(\cU)$. Then 
\bes 
\begin{split}
& DW^1(0, 0, \nu', V_S', \g_{\cC}', V_{\cC}' ) = 
\begin{cases}
- \g'_{\cC} \\
 \p_t (\g_{\cC}') - \chi(\a')\\
 \nu' - \chi(\a')\\
 V_S' - V_{\cC}'\\
 (V_S)^a)' - \chi^a((V_S^1)', \nu_S', \g_{\cC}') \\
 \end{cases}
\end{split}
\ees

  Again, all data on the right are evaluated at $\Si$. Using the triangular structure, it is easy to see this map is surjective. First $\g_{\cC}'$ may be prescribed 
arbitrarily. Next, $H'_{\cC}$, which gives $\a'$ via \eqref{vel}, may be arbitrarily prescribed, as is the trace-free part of $\p_0 \g_{\cC}'$. Having thus fixed the 
choice of $\a'$, $\nu'$ may be arbitrarily prescribed, and similarly for the terms $(V_S^a)'$. Given these choices, $V'_{\cC}$ 
may then also be freely and independently prescribed. This proves the surjectivity. 


The kernel $Ker DW^1$ is thus a closed subspace of a Hilbert space, and so has a closed complement. It then follows from the implicit function theorem for 
Hilbert manifolds that $\cU_c^1 = W^{-1}({\bf 0})$ is a closed Hilbert submanifold of $\cU$. 

\end{proof}

\end{document}
